\documentclass[
  extrafontsizes,
  oneside,
  14pt
]{memoir}

\usepackage{xunicode}

\usepackage[english]{babel}
\usepackage{lmodern}

\usepackage{upquote}

\usepackage{fontspec}

\usepackage{nopageno}

\usepackage{parskip}
\setlength{\emergencystretch}{3em} % prevent overfull lines

\usepackage[
  paperwidth=297mm, % A4 large side
  paperheight=960mm,
  margin=100mm
]{geometry}

\usepackage{graphics}
\graphicspath{ {./images} }

\usepackage{fontspec}
\setmainfont[
  Mapping=tex-text,
  ItalicFeatures={Alternate = 0},
  Ligatures={Common,Rare,Historic},
  % Alternate=1
]{Cormorant Garamond}

\title{\Huge{Départs}}
\author{Jérémie Wenger}
\date{2021}

\begin{document}

\maketitle

\vfill

\noindent Forward/backward neural writing\\
Accompanying stereophonic web app
\\
\\
Python/TensorFlow/p5.js

\noindent Models: OpenAI's GPT-3 (forward); a French GPT-2 model (774m)
trained from scratch on open-source datasets (backward)

\includegraphics{qr}

\newpage

\subsubsection{voix 1}\label{voix-1}

Je partis de Lausanne, pour ne jamais revenir. Tu ne sauras pas où je me suis
rendue. Je ne te dirai pas non plus dans quelle ville je suis morte, il y a de
cela quelques années. Peut-être me chercheras-tu en vain quelque part au bord
de l'océan Atlantique, dans quelque anfractuosité d'une falaise de Bretagne, au
fond d'un tunnel alpin, ou à même la glace aveuglante d'une plage de
Norvège\ldots{} T'obstineras-tu~? Tenteras-tu ta chance sur telle île des
Caraïbes où je m'étais établie quelque temps après avoir quitté Londres, ou
dans l'hôtel miteux du port terrestre de Khorgas, à la frontière sino-kazakhe,
où m'avait entraînée un de ces vieux rêves industriels qui m'avait tant hantée
avant mes années d'exil écossaises~? Une chose est sûre, tu ne me retrouveras
pas.

Le seul conseil que je peux te donner, c'est de ne rien faire, sinon de te
laisser porter par le cours des choses. Quand tu entendras ces paroles, certes,
un grand nombre de choses auront changé, et ce pas forcément en bien. Le monde,
j'en suis sûre, aura pris un tour qui ne te plaira pas. Qui sait, le monde aura
peut-être enfin atteint la crise à laquelle il avait été promis depuis si
longtemps. Le connaissant, il aura même aimé être reconquis par les démons.
Alors, tu pourras t'en vouloir, autant que tu le voudras, de ne pas avoir eu le
courage de céder à la tentation, de ne pas avoir été plus fanatique dans ta
carrière, de ne pas avoir plus légèrement, plus entièrement~!, dédié ton corps
et ton âme à la destruction de ce monde figé auquel tu n'as jamais cru.

Mais alors même que la crise s'aggravera, et risquera de faire basculer les
sociétés industrielles dans ce que les historiens du XXII\textsuperscript{ème}
siècle appelleront «~l'Âge des Turbulences~» -- à moins que, les uns à l'instar
des autres, ils n'adoptent un terme convenu, genre «~planétisation du chaos~»,
pour désigner le flux naissant des siècles de répression et d'obscurantisme
légal qui nous attendent, je ne m'inquiète pas. Je ne m'inquiète plus de grand
chose, il est vrai, surtout depuis la perte de mes malheureux amis, tu sais
bien lesquels, dont la chute, à la fin de mes années urbaines, m'a causé tant
de tristesse. Mais je m'égare. Ce que je veux te dire, c'est que bien des
tentatives d'agir continueront d'échouer~: les tiennes, les miennes, les leurs,
et je n'ai jamais été assez naïve pour croire que l'actif accru de quelques
uns, et le crédit croissant de tous les autres, auraient pu changer quoi que ce
soit au trou béant laissé par les déterminations de l'Histoire, ou suffiraient
à sauver le monde de lui-même.

Comme par le passé, tu n'auras pas à choisir de quel côté de la guerre à venir
tu te rangeras -- comme jadis dans nos disputes --, ni de quel côté se penchera
l'intérêt de ton cœur, comme s'il fallait toujours qu'il y eût un côté.

Pour sûr, si tu suis mes conseils, si tu guides tes pas sur la voie que je t'ai
patiemment tracée, et pour autant que tu ne t'enferres pas dans l'impasse de
poursuivre mon fantôme, tu n'auras pas à choisir. Tu accepteras enfin ce que tu
n'as jamais voulu faire. Tu devras vivre avec tes propres idées, qui ne seront
pas les miennes, même si ce sont elles qui ont créé le berceau qui les a vu
naître. Tu connaîtras le devenir.

Tu te limiteras, tu y seras obligé, à maintenir une modeste neutralité, à
t'abriter le plus longtemps possible dans le contrepoint rassurant de
lamentations que les populations de ton esprit seront plus contraintes que
jamais de gémir, l'échine pliée, les épaules rompues, l'esprit automatique. Te
le reprocheras-tu sans plus finir~? Ou se seront-elles révoltées, par miracle,
ces rampantes créatures de tes imaginaires totalitaires~? Je ne peux le savoir,
même si je ne puis te garantir que personne ne se réveillera, en fin de compte,
un jour ou l'autre, de l'engourdissement de ton monde psychique.

Si je ne me trompe, tout ce que j'ai dit dans ce message n'est à la fin des
fins qu'un leurre destiné à te convaincre de détacher tes yeux des événements
qui se joueront et rejoueront sous tes yeux, pour te tourner vers l'intérieur.
J'anticipe que, comme toujours, tu croiras rencontrer telle femme, avec qui tu
rêveras d'assouvir tel fantasme. J'approuve, tu me connais, et pourtant, c'est
plus fort que moi, l'ennui me gagne. C'est dans ces moments que je te vois
retourner à ta vie déjà bien établie, ou en bonne voie de l'être, à Paris, et
reprendre tes vieilles habitudes.

Promets-moi juste une chose. Alors même que tu te sentiras gagné par la
résignation et la dépression -- ne me contredis pas, je t'en prie, je te sais
l'esclave de ces sirènes --, fais-moi confiance, suis mon conseil, ne t'en fais
pas, laisse-toi aller. Car ce n'est que dans ces moments-là, caressant d'un
doigt moribond le fond de l'abîme, que tu auras une chance de tomber sur une
des perles obscures qui brillent dans les tréfonds. Tu repenseras alors
peut-être au bonheur que tu connaissais jadis dans le confort de mon
inconstance, et tu pleureras enfin, après une vie de sécheresse.

\newpage

\subsubsection{voix 2}\label{voix-2}

À l'heure du réveil, le soleil riait sans doute là-haut dans les pâturages,
trop loin de nous pour qu'on y rêve. Nous, on était pris dans un violent
courant d'air sulfuré, et je sentis dans mes os la même terreur indicible qu'à
chaque fois, une irrépressible obscurité dans l'âme, lorsqu'un avion de chasse
hurla juste au-dessus de nos têtes, condamnant la vallée voisine aux flammes
d'un enfer que nous ne connaissions que trop bien.

«~Supprimons par la force ces infâmes ordures qui menacent de molester la
Pierre d'Unspunnen~!~», tels avaient été les mots du Commandant de Corps à
notre multitude, insouciant magma de troufions immobiles couvrant les plaines
du Mittelland -- cadavres, déjà, presque tous, même si peu d'entre nous le
savaient. Lorsque \emph{Le Temps} avait publié, malgré la réticence de la
rédaction, et uniquement après des menaces répétées de démission en masse de la
part des journalistes, des documents de l'armée contenant des plans
d'éradication émanant directement du Conseil Fédéral, et adressés au général
lui-même, «~à faire exécuter manuellement par vos soins, en rétablissant pour
son implémentation les anciens moyens de l'art abattoire~», n'hésitant pas à
préciser «~qu'il ne fallait pas craindre les objections tracassières des
cantons, ni les embûches fiscales que rencontre tout patriote dans la diligente
mise en place de ces mesures~», on avait poussé des cris d'orfraie comme il se
doit, mais personne ne prit l'affaire réellement au sérieux.

Et alors même que la crise s'aggravait, et risquait de faire basculer les
sociétés industrielles dans ce que les historiens du XXII\textsuperscript{ème}
siècle appelleraient «~l'Âge des Turbulences~» -- à moins que, les uns à
l'instar des autres, ils n'adoptent un terme convenu, genre «~planétisation du
chaos~», pour désigner le flux naissant des siècles de répression et
d'obscurantisme légal qui nous attendent, je ne m'inquiétai pas. Je ne
m'inquiétais plus de grand chose, il est vrai. Par le passé, j'avais craint
d'être impur et ridicule. Pire, j'avais été timide devant le qu'en-dira-t-on. À
ma grande surprise, je ne devais pas le rester. J'avais tenté de briser ces
chaînes par la seule force de ma volonté, en vain. J'avais connu ensuite le
goût fade de la résignation, et m'étais persuadé que ce serait mon régime
minceur pour toujours. Pourtant, depuis que l'idée de la haine s'était
manifestée sur la scène du monde, et que notre pays, comme les autres
d'ailleurs, n'avait plus pour seule passion, pour seule fièvre, que de trouver
un moyen de passer la frontière interdite au-delà de laquelle il n'y a que
l'utopie, je ne rougissais plus.

Je voyais trop où on en était pour ne pas repasser en revue mentalement, comme
un mantra protecteur, les étapes de l'architecture hégélienne. Mon cœur faisait
murmurer à ma bouche ces mots si gris, si abstraits, si intimes, pendant que
mon esprit faisait verser à mes yeux des larmes d'espoir mêlées de sang.

Mécaniste de cœur, navré d'abord de mon étrange connivence avec les
Organicistes de ma section, je ne pouvais m'empêcher de ressentir de la pitié
pour mes compagnons de fortune, que je savais tout comme moi condamnés au
charnier. Je réalisais lentement, malgré les bombes et le brasier, que j'étais
devenu ami avec ceux que malgré tout je viendrais à appeler mes ancêtres.

Car bien que je les aie aimés ardemment, je n'aurais pu dire avoir été lié
d'amitié avec eux jusque-là. Au contraire, c'était dans notre inimité même,
sans nul doute le carburant de notre fascination mutuelle, que j'avais trouvé
mon bonheur durant les premières années de mobilisation.

En l'an 2000, quand tout bascula, j'avais vingt-deux ans. Mon seul désir,
alors, était parvenir à ce niveau supérieur d'effort, requis, me semblait-il,
pour m'expliquer, pour me connaître, pour avoir foi en mes tentatives, pour
vivre~! Je m'engageai dans cette entreprise sans arrière-pensées, obéissant,
pour la première fois depuis l'émergence de ma conscience, plus de quinze ans
auparavant, à toutes les dispositions qu'éveillait en moi la danse
contre-nature de la honte et de la fierté.

Ma vie avait été lamentable, et mes buts étaient restés obstinément
inenvisageables. La guerre n'était, somme toute, qu'une mésaventure de plus.
Pourtant, pendant ce temps de troubles, je compris quelque chose à l'égard de
l'homme que j'allais devenir -- ce double énervant en qui j'allais devoir, bon
gré mal gré, transmigrer.

L'attitude de mes parents vis-à-vis de moi, longtemps insondable,
m'apparaissait maintenant limpide. Trop souvent, dans mon égoïsme et mon
amertume, j'avais épuisé ma mère, déjà assez humiliée par les affres du
travail et l'ignorance hargneuse des hommes. Même la façon dont je
m'adorais moi-même, alors un obstacle insurmontable, semblait s'être
purgée d'elle-même, et je remarquais avec étonnement, chaque fois que
j'étais amoureux de telle femme, avec qui je rêvais d'assouvir tel
fantasme, l'engouement avec lequel j'accueillais la grande influence
qu'elle allait exercer sur ma vie personnelle. J'avais songé déjà à tout
cela au sortir du gymnase, me demandant de plus en plus ce qui m'avait
empêché de me lier avec quiconque, sans que rien ne change. Et tout à
coup, tout était devenu clair, comme si ces questions n'avaient même pas
existé.

J'en éprouvai presque plus que du bonheur~: c'était de la joie, une joie
violente, souvent incontrôlable. J'étais consumé par une soif de tous les
instants. La situation était grave. J'en étais à boire du lait toutes les cinq
heures, et je ne pouvais réprimer des envies de soupes chinoises instantanées.
Que faire~? Ma décision fut aussi soudaine qu'irrévocable. Je sus que c'étaient
eux, mes ancêtres d'élection, toujours avec moi en pensée, qui me donnèrent la
force nécessaire. Je partis de Lausanne, pour ne jamais revenir.

\end{document}
